\subsection{Indecomposable injective modules}

Let A be a ring. The relation “I is a class of indecomposable injective A-modules” is collectivizing (A, X, p. 21, cor. 1); we shall denote by \(\mathcal{I}(A)\) the set of classes of indecomposable injective A-modules.

PROPOSITION 1.--- Let A be a Noetherian ring. For every prime ideal p of A, let \(e_{\mathfrak{p}}: \mathrm{A}/\mathfrak{p} \to \mathrm{I}(\mathfrak{p})\) be an injective envelope of \(\mathrm{A}/\mathfrak{p}\) (A, X, § 1, no. 9).

\begin{enumerate}
\def\labelenumi{\alph{enumi})}
\item
  The A-modules I(p) are indecomposable.
\item
  Let I be an indecomposable injective A-module; the set Ass(I) is reduced to a single element.
\item
  The map \(\mathfrak{p} \mapsto \operatorname{cl}(\operatorname{I}(\mathfrak{p}))\) is a bijection of \(\operatorname{Spec}(\operatorname{A})\) onto \(\mathcal{I}(A)\). The inverse bijection associates to an element I of \(\mathcal{I}(\operatorname{A})\) the unique element of \(\operatorname{Ass}(\operatorname{I})\).
\end{enumerate}

Let \(\mathfrak{p} \in \operatorname{Spec}(A)\); let us prove that the module \(I(\mathfrak{p})\) is indecomposable. By A, X, p. 21, cor. 2, it suffices to prove that if \(\mathfrak{a}\) and \(\mathfrak{b}\) are ideals of A containing \(\mathfrak{p}\) and distinct from \(\mathfrak{p}\), the ideal \(\mathfrak{a} \cap \mathfrak{b}\) is distinct from \(\mathfrak{p}\) ; or if \(a\) is an element of \(\mathfrak{a} - \mathfrak{p}\) and \(b\) is an element of \(\mathfrak{b} - \mathfrak{p}\) the product \(ab\) belongs to \((\mathfrak{a} \cap \mathfrak{b}) - \mathfrak{p}\).

Let I be an indecomposable injective A-module, and let p, q be elements of Ass(I). Then I contains a submodule M isomorphic to A/p and a submodule N isomorphic to A/q. We have M ∩ N ≠ 0 (A, X, p. 21, prop. 14); for every nonzero element x of M ∩ N, we have p = Ann x = q. Since Ass(I) is not empty (IV, § 1, n° 1, cor. 1 of prop. 2), it is reduced to a single element p(I).

We have thus defined two maps. \(\mathfrak{p} \mapsto \operatorname{cl}(I(\mathfrak{p}))\) of \(\operatorname{Spec}(A)\) in \(\mathcal{I}(A)\) and \(I \mapsto \mathfrak{p}(I)\) of \(\mathcal{I}(A)\) in \(\operatorname{Spec}(A)\); let us prove that these two maps are inverse bijections of each other. Let \(\mathfrak{p} \in \operatorname{Spec}(A)\) ; then \(\mathfrak{p}\) belongs to \(\operatorname{Ass}(I(\mathfrak{p}))\), and it is therefore the unique element of \(\operatorname{Ass}(I(\mathfrak{p}))\). Let \(I\) be an indecomposable injective A-module, and \(\mathfrak{p}\) the unique element of \(\operatorname{Ass}(I)\) ; then \(I\) is an injective envelope of \(A/\mathfrak{p}\) (A, X, p. 21, prop. 14). This completes the proof of the proposition.

Remark 1.—Let I be an indecomposable injective A-module, and let p be the unique element of Ass(I); by Prop. 1, I contains a submodule isomorphic to A/p of which it is an injective envelope. In general such a submodule is not unique, as one sees by taking A = Z, p = 0, I = Q.

For every prime ideal \(\mathfrak{p}\) of \(\mathrm{A}\), let us choose, as above, an injective envelope \((\mathrm{I}(\mathfrak{p}), e_{\mathfrak{p}})\) of \(\Lambda/\mathfrak{p}\). By \(\mathrm{A}\), X, p. 22, th. 3, we have:

THEOREM 1.--- Let A be a Noetherian ring. For every injective A-module I, there exists one and only one family of cardinals \((a_{\mathfrak{p}})_{\mathfrak{p}\in\operatorname{Spec}(\mathrm{A})}\), such that I is isomorphic to \(\bigoplus_{\mathfrak{p}}\mathrm{I}(\mathfrak{p})^{(a_{\mathfrak{p}})}\).

By IV, § 1, no. 1, cor. 1 of prop. 3, the set Ass(I) is then the support of the family \((a_{\mathfrak{p}})\) .

Remark 2. Let \(A\) be a Noetherian ring, M an A-module, \(e: \mathrm{M} \to \mathrm{I}\) an injective envelope of M. The set Ass(M) is equal to Ass(I): indeed, the inclusion Ass(M) ⊂ Ass(I) is evident. On the other hand, if \(\mathfrak{p}\) is an element of Ass(I), I contains a submodule N isomorphic to \(A/\mathfrak{p}\) ; since the \(A\)-module \(e^{-1}(\mathrm{N})\) is nonzero, we have \(\operatorname{Ass}(e^{-1}(\mathrm{N})) = \{\mathfrak{p}\}\). This proves that \(\mathfrak{p}\) is associated with \(e^{-1}(\mathrm{N})\), hence to M, whence our assertion.

Let us adopt the notation of the theorem, and suppose moreover that Ass(M) is finite. For every \(\mathfrak{q} \in \operatorname{Ass}(M)\) let us denote by \(Q_{\mathfrak{q}}\) the intersection with M of \(\bigoplus_{\mathfrak{p} \in \operatorname{Ass}(M) - \{\mathfrak{q}\}} I(\mathfrak{p})^{(a_{\mathfrak{p}})}\). Then \((Q_{\mathfrak{q}})_{\mathfrak{q} \in \operatorname{Ass}(M)}\) is a reduced primary decomposition of 0 in M (IV, § 2, n° 3, def. 3). Indeed, we have \(\cap Q_{\mathfrak{q}} = 0\); as \(M/Q_{\mathfrak{q}}\) is identified with a nonzero submodule of \(I(\mathfrak{q})^{(a_{\mathfrak{q}})}\), one has \(\operatorname{Ass}(M/Q_{\mathfrak{q}}) = \{\mathfrak{q}\}\) and it suffices to apply prop. 4 of loc. cit.

Example.--- Let A be a principal ideal ring, and let K be its field of fractions. The injective A-modules are the divisible A-modules (A, X, p. 17, cor. 2). The A-module K is an injective envelope of A (A, X, p. 20, example 1). Let \(p\) be an extremal element of A, and \(\mathfrak{p}\) the (maximal) ideal \(Ap\). Let us denote by \(e: A/\mathfrak{p} \longrightarrow K/A_{\mathfrak{p}}\) the homomorphism that maps the class of an element \(a \in A\) to the class of \(a/p\). Let us prove that \((K/A_{\mathfrak{p}}, e)\) is an injective envelope of \(A/\mathfrak{p}\). The homomorphism \(e\) is injective. The A-module \(K/A_{\mathfrak{p}}\) is a quotient of a divisible module, hence is divisible. Let \(x\) be a nonzero element of \(K/A_{\mathfrak{p}}\) ; it is the class of an element \(a/p^n\) of K, with \(a \in A - \{0\}\) and \(n \geqslant 1\) (A, VII, p. 10, th. 2). We then have \(p^{n-1}x = e(a)\) , hence \(e(Ax) \neq 0\), which proves our assertion.

It then follows from th. 1 that every divisible A-module is a direct sum of A-modules isomorphic to K or to \(K/A_{p}\) for a maximal (principal) ideal p of A.

